\documentclass[11pt]{article}
\usepackage{amsmath,amssymb}
\title{How many histories can rule out a timed recovery?}
\date{25 November 2026}
\begin{document}
\maketitle
\paragraph{Problem.}
A receiving controller must choose a single recovery mode \(m\in M\)
and a start tick \(t\) from its current view. A nonempty finite class
\(H\) of actual source histories is compatible with that view.
For each mode and history, assume the complete admissible starts
are the integer interval \([L_m(h),U_m(h)]\), possibly empty.
For example, a mode can require physically available support
from its documented readiness until a verified occupancy cutoff,
with a fixed transfer dwell. Earliest readiness gives a lower
endpoint; the cutoff minus dwell is an upper endpoint.
Other physical or lawful requirements must already be reflected
in the endpoints. The class requires a complete single interval
for each fixed mode, not a convex hull of a disconnected feasible
set. Choosing a mode and time is a one-cut decision; this is not
completeness of a later history-dependent policy.

\paragraph{Exact decision and a small rejecting witness.}
In mode \(m\), the common feasible interval is
\([\max_hL_m(h),\min_hU_m(h)]\).
Thus the earliest start in a nonempty intersection is a witness.
If there is no admissible choice, retain a history attaining the
maximum lower endpoint and a history attaining the minimum upper
endpoint in each mode (the same history can supply both).
The at most \(2|M|\) histories together already refute every mode.
The witness has only two named endpoint obligations per mode;
it need not retain all their other fields.
This is the interval Helly argument applied separately to each
allowed alternative, not a new theorem about arbitrary control sets.

The bound is tight even if every proper subfamily is compatible.
For each mode \(m\), create two histories whose windows in that mode
are \([0,0]\) and \([1,1]\) respectively, and give both of them
window \([0,1]\) in every other mode.
All \(2|M|\) histories exclude every choice. Omitting a history
permits the opposite tick in its affected mode.
In particular if there are at least two modes, every pair can have
a feasible recovery even while the full class has none.
This can arise from choosing between alternative pads: during an
unchanged one-tick transfer, one history closes the relevant pad
after tick zero while another does not make it ready until tick one.
No assumption of independent failure of the pad constraints is needed.

\paragraph{What survives a source replacement.}
For a fixed mode set, a rejecting witness survives if its designated
histories are still admitted and the witnessed lower endpoints
do not move earlier and the witnessed upper endpoints do not move
later. Other rows can change, and more histories can be admitted
without restoring the rejected mode. The same check is insufficient
if a witness has been removed from the information class or a new
mode is authorized. Losing the old witness does not prove existence
of a recovery, and is not a reason to accept without revalidation.
Validity of these endpoint facts still requires a source mechanism;
a self-declared update flag is not evidence that there are no
other edits or that an old history remains relevant.

\paragraph{Index and comparison.}
With an independently established complete stream of known endpoint
updates, keep two addressed heaps per mode, one for each extremum.
Replacing one window takes \(O(\log|H|)\) comparisons per heap;
scanning at most \(|M|\) roots decides admission or emits the witness.
The heaps hold \(O(|M||H|)\) current entries and do not accumulate
stale generations. This calculation excludes acquiring, authenticating
or identifying the source update; an unreported sparse edit is a
different discovery problem. An equally factored direct index ties
the bound. Full enumeration serves as an exact diagnostic, not as
the strongest performance comparator.
The checker compares certificates with enumeration on 6,400
instances, includes the tight family for 1 through 16 modes,
and checks 16,384 indexed replacements.
Lean checks exact extrema, the rejecting witness, sharpness,
and the replacement facts. It does not verify the Python heaps,
supply physical endpoint laws, or discharge R1/R2/I1.
\end{document}
